\section*{अध्याय \ref{ch:b3:stable-isotopes} --- समस्थानिक भू-रसायन}

\begin{solution}{exo:b3:stable-isotopes:1}
$\alpha = 1030.5/997.9 = 1.032669$; $1000\ln\alpha = 32.15$; $\Delta = 30.5 -
(-2.1) = \qty{32.6}{\permil}$। ये \qty{0.45}{\permil} से भिन्न हैं, जो
प्रतिज्ञप्ति के $(\delta_A^2 - \delta_B^2)/2000 = 0.46$ के बराबर है।
\end{solution}

\begin{solution}{exo:b3:stable-isotopes:2}
$\delta^{18}\mathrm O = 0.69\times(-25) - 13.6 = \qty{-30.9}{\permil}$;
$\delta\mathrm D = 5.6\times(-25) - 100 = \qty{-240}{\permil}$; $d = -240 +
8\times30.85 = \qty{6.8}{\permil}$।
\end{solution}

\begin{solution}{exo:b3:stable-isotopes:3}
$\Delta_B = -26.0 - 1.0 = \qty{-27.0}{\permil}$; $f_{\mathrm{org}} = (-5 -
1.0)/(-27.0) = 0.22$।
\end{solution}

\begin{solution}{exo:b3:stable-isotopes:4}
$T^2 = 0.0391\times10^6/(0.600 - 0.154) = \num{87668}$, $T = \qty{296.1}{K}$,
\qty{22.9}{\celsius}।
\end{solution}

\begin{solution}{exo:b3:stable-isotopes:5}
$f^{0.01} = 967/987$, $f = \exp(100\ln0.97974) = 0.129$: \qty{87}{\percent} वाष्प
वर्षित हो चुका है। वर्षा: $1.0100\times967 - 1000 =
\qty{-23.3}{\permil}$।
\end{solution}

\begin{solution}{exo:b3:stable-isotopes:6}
$\alpha = 1000/960 = 1.041667$; $1000\ln\alpha = 40.82$; $1000(\alpha - 1) =
41.67$; $\Delta = 40$। दोनों सन्निकटन लगभग \qty{0.8}{\permil} से चूकते हैं,
दूसरा पहले से दोगुना चूकता है: \qty{40}{\permil} पर प्रथम कोटि तक का
प्रसार अब पर्याप्त नहीं रहता।
\end{solution}

\begin{solution}{exo:b3:stable-isotopes:7}
$\delta_c = (1.50 + 29.99)/0.97001 = \qty{32.46}{\permil}$ VSMOW; $1000\ln\alpha
= 1000\ln1.032464 = 31.95$; $T = 18030/(31.95 + 32.42) = \qty{280.1}{K}$,
\qty{7.0}{\celsius}। प्रभाजन प्रति केल्विन लगभग \qty{0.23}{\permil} बदलता है, अतः
जल पर \qty{0.1}{\permil} का परिवर्तन लगभग \qty{0.4}{\celsius} के बराबर है।
\end{solution}

\begin{solution}{exo:b3:stable-isotopes:8}
हिम-आयतन $125\times0.011 = \qty{1.375}{\permil}$; शेष, $1.77 - 1.375 =
\qty{0.40}{\permil}$। \qty{2}{\celsius} के निकट ढलान $18030/275.15^2 =
\qty{0.238}{\permil.K^{-1}}$ है, इसलिए गहरा जल लगभग \qty{1.7}{\celsius}
ठंडा था।
\end{solution}

\begin{solution}{exo:b3:stable-isotopes:9}
$\Delta^{33}\mathrm S = 8.2 - 0.515\times10.0 = \qty{3.05}{\permil}$, 0 से दूर:
द्रव्यमान-स्वतंत्र सल्फ़र, जो केवल ओज़ोन परदे के बिना बनता है। पाइराइट
ऐसे वायुमंडल में बना जिसमें पर्याप्त ऑक्सीजन नहीं थी, लगभग
\qty{2.45}{\Ga} से पहले।
\end{solution}

\begin{solution}{exo:b3:stable-isotopes:10}
बंद तंत्र में भारी समस्थानिक संरक्षित रहता है: $R_0 = fR_v + (1 -
f)\alpha R_v$, अतः $R_v/R_0 = 1/[f + \alpha(1 - f)]$। $f = 0.2$, $\alpha =
1.01$: $1/1.008$, वाष्प \qty{7.9}{\permil} हल्का है; रैले से
$0.2^{0.01}$, \qty{16.0}{\permil} हल्का। प्रत्येक संघनित को तुरंत हटाने से
वाष्प दोगुनी दूरी तक चला जाता है।
\end{solution}

\begin{solution}{exo:b3:stable-isotopes:11}
$M = m(\delta_s - \Delta)/\Delta = 2000\times(-60 + 2.5)/(-2.5) = \qty{46000}{Gt}$।
फिर $m = 46000\times(-2.5)/(-25 + 2.5) = \qty{5100}{Gt}$, जो प्रकाशित
6000 से \qty{12000}{Gt} तक की सीमा से थोड़ा नीचे है, जिसकी ऊपरी सीमा बड़े
विचलन से मेल खाती है।
\end{solution}

\begin{solution}{exo:b3:stable-isotopes:12}
$\delta_v = (1000 + \delta_{v,0})\exp\int(\alpha - 1)\,\mathrm d\ln f - 1000$, अतः
$\mathrm d\delta_v = (1000 + \delta_v)(\alpha - 1)\,\mathrm d\ln f$, और $f =
e_s/e_{s,0}$ से $\mathrm d\ln f = \mathrm d\ln e_s$। \qty{-20}{\celsius} पर,
$\alpha = 1.01872$: $0.01872\times950.9\times0.0960 =
\qty{1.71}{\permil.K^{-1}}$ वाष्प के लिए।
\end{solution}

\begin{solution}{pb:b3:stable-isotopes:1}
\textbf{1.} $\alpha^{18}$: 1.01605, 1.01691, 1.01779, 1.01872, 1.01968,
1.02068; $\alpha^{\mathrm D}$: 1.1424, 1.1524, 1.1630, 1.1744, 1.1866, 1.1997
($-5$ से \qty{-30}{\celsius} तक)।

\textbf{2.} $980\times0.425^{0.01605} - 1000 = \qty{-33.4}{\permil}$।

\textbf{3.} $1.01691\times966.6 - 1000 = \qty{-17.0}{\permil}$।

\textbf{4.} वाष्प $966.6\times(0.3976)^{0.01779} - 1000 = \qty{-49.1}{\permil}$;
हिम $1.01872\times950.9 - 1000 = \qty{-31.3}{\permil}$।

\textbf{5.} वाष्प $950.9\times(0.3680)^{0.01968} - 1000 = \qty{-67.6}{\permil}$;
हिम $1.02068\times932.4 - 1000 = \qty{-48.3}{\permil}$।

\textbf{6.} $(932.4/980)\times0.0622 = 0.059$: ऑक्सीजन-18 का \qty{5.9}{\percent},
जल के \qty{6.2}{\percent} के विरुद्ध: वाष्प ने
भारी समस्थानिक को अधिक तेज़ी से खोया है।

\textbf{7.} $(-17.0 + 48.3)/20 = \qty{1.57}{\permil}$ प्रति डिग्री।

\textbf{8.} $t = (-48.3 + 13.6)/0.69 = \qty{-50}{\celsius}$। यह संबंध
सतही तापमानों के लिए है; हिमचादर पर सतह उस वायु से बहुत ठंडी होती है
जहाँ हिम बनता है, और मॉडल मिश्रण तथा गतिज प्रभावों की उपेक्षा करता है।

\textbf{9.} वाष्प $-247.5$, $-352.5$, \qty{-462.7}{\permil}; हिम $\delta\mathrm D
= -132.8$, $-239.6$, \qty{-355.4}{\permil}।

\textbf{10.} $d = -132.8 + 8\times17.02 = 3.3$; फिर 10.7 और
\qty{31.3}{\permil}: आधिक्य अंतर्देशीय बढ़ता है। आसवन की स्थानीय ढलान,
$\mathrm d\delta\mathrm D/\mathrm d\delta^{18}\mathrm O = (\alpha^{\mathrm D} -
1)(1000 + \delta\mathrm D)/[(\alpha^{18} - 1)(1000 + \delta^{18}\mathrm O)]$, जो
8 से नीचे गिर जाती है जब वाष्प बहुत हल्की हो जाती है ($1000 + \delta\mathrm D$ कहीं अधिक तेज़ी से
घटता है $1000 + \delta^{18}\mathrm O$ की तुलना में), इसलिए हिम रेखा के ऊपर खिसक जाता है।

\textbf{11.} रेखा $8\times(-48.34) + 10 = \qty{-376.7}{\permil}$ देती है;
हिम उससे \qty{21}{\permil} ऊपर है। मॉडल में केवल साम्य प्रभाजन है;
वास्तविक हिम अतिसंतृप्त वायु में, गतिज प्रभाव के साथ, बनती है,
और वायु राशियाँ रास्ते में दूसरी वायु राशियों से मिश्रित होती हैं।

\textbf{12.} $\delta_c = (3.25 + 29.99)/0.97001 = 34.27$; $1000\ln\alpha =
33.69$; $T = 18030/66.11 = \qty{272.7}{K}$, \qty{-0.4}{\celsius}।

\textbf{13.} $\delta_c = 36.09$, $1000\ln\alpha = 35.46$, $T = \qty{265.6}{K}$,
\qty{-7.5}{\celsius}: यह उस तापमान से ठंडा है जिस पर समुद्र-जल बिना जमे रह सकता है।
इस परिवर्तन को पूरी तरह तापमान के कारण नहीं माना जा सकता।

\textbf{14.} $\delta_w = 125\times0.011 = \qty{+1.375}{\permil}$; $1000\ln(1036.09/1001.375)
= 34.08$; $T = 18030/66.50 = \qty{271.1}{K}$, \qty{-2.0}{\celsius}: तापमान में
\qty{1.6}{\celsius} की कमी।

\textbf{15.} $\Delta_{47} = 0.0391\times10^6/272.71^2 + 0.154 = 0.6797$ और
$0.0391\times10^6/271.12^2 + 0.154 = 0.6859$: दोनों मानों का अंतर
\qty{0.006}{\permil} है।

\textbf{16.} यह मापता है कि भारी समस्थानिक खनिज के भीतर
किस प्रकार व्यवस्थित हैं, यह एक आंतरिक साम्य है जो केवल तापमान पर निर्भर करता है।

\textbf{17.} $f_{\mathrm{org}} = (-5 - 0.5)/(-28) = 0.20$।

\textbf{18.} $m = 46000\times(-3.0)/(-60 + 3.0) = \qty{2400}{Gt}$।

\textbf{19.} $46000\times(-3.0)/(-25 + 3.0) = \qty{6300}{Gt}$।

\textbf{20.} $46000\times(-3.0)/(-5 + 3.0) = \qty{69000}{Gt}$, जो स्वयं भंडार
से भी अधिक है, और वह भी भूवैज्ञानिक रूप से संक्षिप्त समय में: अविश्वसनीय है।

\textbf{21.} क्रोड में कार्बोनेट विलयन की गहराई और विस्तार, जो
मापता है कि कितना कार्बन महासागर में प्रविष्ट हुआ, उसकी संरचना से स्वतंत्र रूप से;
विचलन के आकार के साथ यह $\delta_s$ निर्धारित करता है।

\textbf{22.} डोम इल्से की हिम में $\delta^{18}\mathrm O =
\textbf{\qty[text-series-to-math=true]{-48.3}{\permil}}$ (\qty{-48.4}{\permil}
जबकि $\alpha$ लगातार बदलता है)।
\end{solution}
